\section{模型验证、稳健性与工程交付}

本项目将“数值求解正确”与“提交材料正确”分开审计。数值层面，Q1 通过精确公共坐标匹配和合成偏差恢复验证时间同步；Q2/Q3 通过 NIS、白化创新相关、协方差正定性、敏感性分析和合成数据检验状态融合；Q4 对所有最终两位小数任务时刻重新以 0.01 s 细网格检查完整准备区间，并对模板未写区域做值与样式快照比较。

交付层面新增三类自动检查：
\begin{enumerate}[label=(\arabic*)]
  \item \textbf{输入审计：}将本地附件的 SHA256、字节数、sheet 名、行数与字段和 \texttt{00\_problem/input\_manifest.json} 对照；
  \item \textbf{结果审计：}检查 Q1--Q3 轨迹的有限性、严格递增和 10 Hz 步长，检查统一时间偏差转换、Q4 模板容量和 MILP gap，并确认每个正式图件清单中的 PDF 存在；
  \item \textbf{提交审计：}在打包前生成 JSON/Markdown 报告和 SHA256 文件清单，避免把历史论文、临时文件或错误版本的 \texttt{result.xlsx} 混入最终 ZIP。
\end{enumerate}

论文数值通过 Python 自动生成的 TeX 宏与表格注入，而不是在多个章节手工重复录入。因此当正式结果文件变化时，重新运行资产生成器即可同步论文，降低“代码结果已更新、正文仍保留旧数字”的风险。

需要强调的是，模型内位置标准差和 Q4 扰动试验反映的是给定状态空间模型及压力情景下的稳定性，不应被解释为覆盖全部模型不确定性的严格概率保证。系统偏差结论同样是条件于时间配准结果和预设统计/工程阈值的判定。
